\section{Introduction}

AI-mediated online platforms are beginning to allocate work to LLM agents that act for different principals: internal teams, customers, vendors, analysts, and platform services. These agents do not operate as isolated chatbots. They read and write shared memories, workspace documents, summaries, messages, tool outputs, and final updates. This shared runtime state is the source of their coordination value, but it also changes the form of privacy risk. A sensitive budget, customer reference, delay rationale, or credential-like marker can first enter a shared artifact, then be summarized, forwarded, or used by another agent before a final answer is ever produced.

This paper studies that failure mode as a \emph{privacy propagation externality}. A writer may benefit locally from placing detailed information into a shared object, while the downstream exposure cost is borne by another principal whose information becomes visible to unauthorized agents or external channels. The cost is shaped by the platform topology: the same raw value has different consequences when written to private memory, sent through a one-hop chain, aggregated by a star coordinator, or placed on a blackboard that many agents can read. This makes privacy governance a runtime platform problem rather than only a final-output moderation problem.

Existing agent-safety work has identified indirect prompt injection, tool misuse, memory poisoning, and multi-agent compromise \cite{greshake2023indirect,liu2023promptinjection,debenedetti2024agentdojo,chen2024agentpoison,gu2024agentsmith,zhang2025asb}. However, defenses and evaluations often emphasize whether the final action or final answer is unsafe. Final-surface checks miss cases in which private content has already propagated through intermediate state. Static access control is also incomplete: a channel label alone does not reveal the payload's provenance, the recipient's role, the fanout of the target object, or whether policy permits only an abstraction rather than the raw value.

We present \textbf{FlowFence}, a runtime mediation mechanism for AI-mediated platforms. FlowFence intercepts information-moving events before they update shared state or reach external channels. For each event, it combines policy checks with provenance, topology, privilege, fanout, and semantic-request signals. The guard can allow the event, rewrite it into a validated safe view, quarantine raw content, block the transition, or narrow its propagation rights. The key design choice is to govern the movement of information, not merely the final text: a vendor-facing update may safely state that a budget constraint exists while withholding the exact internal budget.

We evaluate FlowFence on retrieval-memory and multi-agent workflow settings. The main workflow models an enterprise vendor-update task with chain, star, and blackboard communication structures. MiniMax is used only as a final writer; guard decisions and leakage metrics are computed by deterministic runtime code. FlowFence preserves deterministic task success and eliminates unauthorized raw and external leakage in the evaluated FlowFence runs. A topology slice under no defense shows that broader blackboard-style sharing produces higher raw leakage and cascade size than chain and star structures. A label-free validation further shows that containment does not depend on evaluator-only attack annotations.

The paper makes three contributions:
\begin{itemize}
  \item We formulate privacy leakage in multi-principal LLM platforms as event-level propagation over runtime graphs, capturing how locally useful writes can impose downstream exposure costs on other principals.
  \item We develop FlowFence, a runtime mediation mechanism that combines policy-aware safe-view construction, quarantine, propagation-right narrowing, and topology-aware risk scoring.
  \item We provide empirical evidence across retrieval-memory and multi-agent workflow benchmarks, including a topology analysis and label-free validation, showing that intermediate propagation can be contained while preserving the vendor-facing task.
\end{itemize}
